\subsection{Extend a known finite release channel}
\label{sec:commitment}
Suppose the pooled earlier output $O$ has known conditional law
$P(o\mid s)$. It includes disclosures to every covered recipient;
the finite construction can treat this bundle as one old channel.
A new workload requires $r$ independent uniform binary coordinates
$Z_1,\ldots,Z_r$, independent of the old protected input $s$.
The agency may access the data again, but must retain the actual old
output. The complete pair of old output and new decisions must satisfy
pure $E$-DP under replacement of $(s,Z)$. We optimize minimum new-task
accuracy, averaging over a declared public prior on $s$.

This construction treats $(s,Z)$ as one protected unit and compares
every pair of its values. Exposing its separately randomized row outputs
gives the local-DP observation model, even when the agency samples them.
It answers how to extend an irrevocable finite disclosure, rather than
which mechanism the agency should choose before any disclosure exists.
For population learning, a central computation that never exposes
row-level responses is a necessary comparator. The model-history
formulation below handles a different, general record-neighbor graph;
its extension principle does not require a local old release.

Put $N=2^r$ and $\lambda=e^E$. Symmetry groups new error vectors by
their number $h$ of incorrect coordinates. For each old output choose
a probability floor
\[
 \frac{\max_sP(o\mid s)}{N\lambda}
 \le m_o\le \frac{\min_sP(o\mid s)}N.
\]
For each old state, assign that floor to every error vector and fill
remaining probability in increasing order of $h$, up to $\lambda m_o$
per vector. The optimum occurs at an endpoint or one of finitely many
scalar breakpoints. Proposition~\ref{prop:finite-commitment} gives the
formula, boundary cases and proof. For binary randomized response (RR)
as the old channel, the calculation takes $O(r^2)$ arithmetic operations.
It requires the complete old channel, not merely its privacy scalar.
It is not a tractable characterization of arbitrary trained weights.

The one-new-coordinate formula recovers prior work
\citep{yamamoto2024privacyoptimized}. For several coordinates the
proper comparison is an optimized joint mechanism for the new vector,
independent of the old response, not only coordinate-wise RR.
\subsection{Use the actual retained model history}
\label{sec:weights-main}
Let $P_D(m)$ be the joint law of the pooled old history, including
weights and other data-dependent disclosures to any covered recipient.
Separate models' marginal laws do not determine this joint law.
An extension $Q_D(m,a)$
must satisfy
\[
 \begin{aligned}
 \sum_a Q_D(m,a)&=P_D(m),\\
 Q_D(m,a)&\le e^E Q_{D'}(m,a)\quad(D\sim D').
 \end{aligned}
\]
Given the actual $m$, the agency can realize a feasible finite extension
by sampling $a$ from $Q_D(m,a)/P_D(m)$ on its positive support. This
requires protected data access and knowledge of the conditional laws;
it is not free post-processing of $m$.

Summing the joint inequality over $a$ shows that the old law must itself
satisfy the requested bound. A continuation cannot repair a nonprivate
earlier release. For example, a deterministic old output that differs on
neighbors already gives an event of probability one versus zero and
cannot admit a finite pure-DP extension. Knowing a law is therefore
necessary but not sufficient for feasibility.

Protecting a full private training representation is sufficient but can
be unnecessarily restrictive when only particular weights were exported
and no recipient obtained that representation.
If private training data releases are shared, they remain part of the old
disclosure history; a model-only comparison cannot omit them.
Appendix~\ref{sec:actual-model-history} proves a strict separation in a
three-record learning problem and a recoverability condition for equality.
An omitted disclosure invalidates the comparison. The result is not a
budget refund or a general model analyzer; ordinary composition remains
a conservative baseline.

For a fixed $q$-value record alphabet, suppose the old channel and the
utility objective are invariant to record permutations. Averaging a
feasible extension over those permutations preserves its old marginal,
record-level privacy and utility. The channel can then depend only on
the $q$ category counts, with $\binom{n+q-1}{q-1}$ states; a neighboring
state moves one record between categories. Proposition~\ref{prop:count-reduction}
proves the reduction and gives the necessary multiplicity weights.
For our two-bit records and two binary model outputs, this gives
$4\binom{n+3}{3}$ variables instead of $4^{n+1}$.
It requires invariance of the complete earlier disclosure, not merely
an informal claim that a model ignores row order. Count aggregation
removes row-order redundancy, but still covers every possible count
state, rather than just the observed count vector. Both the record count
$n$ and alphabet size $q$ limit tractability.

With uncertain public utility laws, optimize a mixture or worst-case
regret over a declared set; the privacy constraints stay unchanged.
Choosing a task from protected data instead needs a separate charge.
Appendix~\ref{app:uncertain-workloads} proves the bound for a fixed menu
of joint channels and a private selector after the old export.

